\documentclass[11pt]{letter}
\usepackage[margin=1in]{geometry}
\usepackage[hidelinks]{hyperref}
\usepackage{url}
\setlength{\parskip}{0.85em}
\setlength{\parindent}{0pt}
\begin{document}

Dear Editors,

Please consider our manuscript, ``Four-Corner Spectral Tuning for Quadratic ADMM with Application to Diffeomorphic Image Registration,'' for publication as a Regular Paper in the \emph{Journal of Mathematical Imaging and Vision}.

Selecting the penalty and relaxation parameters of over-relaxed ADMM (oADMM) can require repeated spectral-radius evaluation of a large iteration operator. For quadratic splits whose attained joint modal spectrum is a Cartesian product of data and regularizer curvatures, we show that the worst modal contraction is attained among only four endpoint combinations. Analytic relaxation and a finite algebraic candidate set then replace grid-based or iterative penalty search.

Diffeomorphic image registration provides the motivating noncommuting application: the assembled optical-flow Hessian varies spatially, while the Sobolev regularizer couples neighboring pixels. The exact four-corner theorem therefore does not apply to the full operator. A family of globally coefficient-frozen surrogates indexed by the local optical-flow Hessian blocks, together with the rank-one data Hessian, nevertheless yields an inexpensive one-shot predictor determined by four endpoint curvatures.

On all 134 FIRE retinal pairs at $256\times256$, one-shot four-corner tuning reduced median pairwise runtime by 34.4\% and inner iterations by 39.6\% relative to a validation-selected fixed comparator chosen on separate images, while landmark accuracy was essentially unchanged and no nonpositive discrete Jacobian was observed. On ten CIMA histology pairs, matched-protocol runtime decreased by 20.1\%. Supplementary Information contains certificate constructions, implementation details, and additional diagnostics.

We believe the manuscript is well suited to JMIV. It combines a structural spectral analysis of operator splitting with a concrete variational imaging application and reproducible evaluation on public registration benchmarks. The approach is complementary to general numerical spectral ADMM optimization: stronger modal assumptions yield an exact four-corner reduction and a practical predictor for spatially varying registration.

Source code, processed outputs, and preparation and evaluation scripts are available at \url{https://github.com/xhxuciedu/admm}.

The authors confirm that this work is original, unpublished, and not under consideration elsewhere, and that it does not extend a prior conference paper. There is no specific funding and no competing interest.

Thank you for considering our manuscript.

Sincerely,\\[1.25em]
Xiaohui Xie\\
Department of Computer Science\\
University of California, Irvine\\
Irvine, California, USA\\
\href{mailto:xhx@uci.edu}{xhx@uci.edu}

\end{document}
